\documentclass[aspectratio=169,10pt]{beamer}
\usetheme[numbering=fraction,progressbar=frametitle]{metropolis}
\usepackage{fontspec}
\setsansfont{Latin Modern Sans}
\usepackage{amsmath,amssymb,mathtools,graphicx,booktabs}
\usepackage{polyglossia}
\setdefaultlanguage{french}
\setbeamertemplate{navigation symbols}{}
\setbeamerfont{frametitle}{size=\large}
\setbeamerfont{normal text}{size=\normalsize}
\setlength{\parskip}{.45em}
\hypersetup{pdfauthor={},pdfsubject={Lecture slides},pdfkeywords={},pdftitle={Limites et continuité}}
\providecommand{\N}{\mathbb{N}}
\providecommand{\Z}{\mathbb{Z}}
\providecommand{\D}{\mathbb{D}}
\providecommand{\Q}{\mathbb{Q}}
\providecommand{\R}{\mathbb{R}}
\providecommand{\C}{\mathbb{C}}
\providecommand{\vect}[1]{\overrightarrow{#1}}
\providecommand{\vi}{\vec{\imath}}
\providecommand{\vj}{\vec{\jmath}}
\providecommand{\vk}{\vec{k}}
\providecommand{\norm}[1]{\left\lVert #1\right\rVert}
\providecommand{\vc}[2]{\begin{pmatrix}#1\\#2\end{pmatrix}}
\providecommand{\deter}[2]{\det\left(#1,#2\right)}
\providecommand{\abs}[1]{\left\lvert #1\right\rvert}
\providecommand{\dd}{\,\mathrm{d}}
\providecommand{\e}{\mathrm{e}}
\providecommand{\Card}{\operatorname{Card}}
\providecommand{\E}{\operatorname{E}}
\providecommand{\V}{\operatorname{V}}
\providecommand{\degre}{^{\circ}}
\providecommand{\un}[1]{\,\mathrm{#1}}
\providecommand{\sci}[2]{#1\times10^{#2}}
\providecommand{\pgcd}{\operatorname{PGCD}}
\providecommand{\ppcm}{\operatorname{PPCM}}
\title{Limites et continuité}
\subtitle{Cours et exemples}
\date{}
\author{}
\begin{document}
\begin{frame}[plain]
\vfill
{\usebeamerfont{title}\inserttitle\par}
\medskip
{\usebeamerfont{subtitle}\insertsubtitle\par}
\medskip
{\usebeamercolor[fg]{progress bar}\rule{\textwidth}{.5pt}}
\vfill
\end{frame}
\begin{frame}{Objectifs}
Calculer une limite en un point ou à l’infini, distinguer les limites latérales et étudier la continuité.
\end{frame}
\begin{frame}{Observer un rapprochement}
Dire que \(f(x)\) tend vers \(\ell\) lorsque \(x\) tend vers \(a\), c’est dire que l’on peut rendre \(f(x)\) aussi proche que l’on veut de \(\ell\) en prenant \(x\) suffisamment proche de \(a\) dans le domaine, avec \(x\ne a\). On note \(\lim_{x\to a}f(x)=\ell\). On suppose que des points du domaine autres que \(a\) peuvent s’approcher de \(a\).
\end{frame}
\begin{frame}{Observer un rapprochement}
La limite ne dépend pas de la valeur éventuelle \(f(a)\). Les limites à gauche (\(x<a\)) et à droite (\(x>a\)) doivent coïncider pour qu’une limite bilatérale existe lorsque les deux côtés appartiennent au domaine.
\end{frame}
\begin{frame}{Observer un rapprochement}
Pour \(f(x)=\dfrac1{x-2}\), \(\lim_{x\to2^-}f(x)=-\infty\) et \(\lim_{x\to2^+}f(x)=+\infty\). Ces limites sont différentes : il n’existe pas de limite bilatérale en \(2\).
\end{frame}
\begin{frame}{Limites à l’infini et opérations}
\(f(x)\to+\infty\) signifie que \(f(x)\) dépasse tout seuil réel fixé lorsque \(x\) est suffisamment proche du point considéré, ou suffisamment grand si \(x\to+\infty\). La définition de \(-\infty\) utilise des valeurs inférieures à tout seuil fixé.
\end{frame}
\begin{frame}{Limites à l’infini et opérations}
Si \(f\to\ell\) et \(g\to m\) avec \(\ell,m\in\mathbb{R}\), alors \(f+g\to\ell+m\), \(fg\to\ell m\), et \(f/g\to\ell/m\) si \(m\ne0\). Si \(f\ge0\) et \(f\to\ell\ge0\), alors \(\sqrt f\to\sqrt\ell\). Les polynômes ont pour limite leur valeur en tout réel. À l’infini, \(x^n\) (\(n\in\mathbb{N}^*\)) domine ses termes de degré inférieur. Pour un quotient de polynômes, factoriser les puissances dominantes.
\end{frame}
\begin{frame}{Limites à l’infini et opérations}
Les écritures \(0/0\), \(\infty/\infty\) et \(\infty-\infty\) sont des formes indéterminées : elles ne donnent aucune réponse. \(\infty\) n’est pas un nombre sur lequel on calcule comme dans \(\mathbb{R}\).
\end{frame}
\begin{frame}{Limites à l’infini et opérations}
Au voisinage d’un réel, factoriser puis simplifier sans oublier la valeur exclue; avec des racines, multiplier par l’expression conjuguée. À l’infini, diviser numérateur et dénominateur par la plus grande puissance utile.
\end{frame}
\begin{frame}{Limites à l’infini et opérations}
Pour \(x\ne4\) et \(x\ge0\),
\[
\dfrac{\sqrt x-2}{x-4}=\dfrac1{\sqrt x+2},\qquad
\lim_{x\to4}\dfrac{\sqrt x-2}{x-4}=\dfrac14.
\]
De même, \(\dfrac{3x^2-x+1}{2x^2+5}\to\dfrac32\) quand \(x\to\pm\infty\), car on divise par \(x^2\).
\end{frame}
\begin{frame}{Continuité}
Une fonction définie en \(a\) est continue en \(a\) lorsque \(\lim_{x\to a}f(x)=f(a)\), la limite étant prise dans son domaine. Elle est continue sur un intervalle si elle est continue en chacun de ses points, avec une limite latérale aux extrémités incluses.
\end{frame}
\begin{frame}{Continuité}
Les fonctions polynômes sont continues sur \(\mathbb{R}\). Les quotients de fonctions continues sont continus là où le dénominateur ne s’annule pas. La racine carrée est continue sur \([0;+\infty[\). Les fonctions \(\sin\) et \(\cos\) sont continues sur \(\mathbb{R}\), et \(\tan\) sur chacun des intervalles de son domaine.
\end{frame}
\begin{frame}{Continuité}
Si \(f\) est continue sur \([a;b]\), tout réel compris entre \(f(a)\) et \(f(b)\) est atteint au moins une fois. Si \(f\) est de plus strictement monotone, cet antécédent est unique.
\end{frame}
\begin{frame}{Continuité}
\begin{enumerate}
\setcounter{enumi}{0}
\item Écrire le domaine et le point approché. Une limite en une borne peut n’être accessible que d’un côté.

\item Si les opérations portent sur des limites finies et un dénominateur non nul, substituer puis citer la règle utilisée.

\item Pour un dénominateur tendant vers zéro, déterminer son signe de chaque côté; une taille ne suffit pas à déterminer le signe d’une limite infinie.

\item Pour une forme indéterminée, transformer l’expression sur un voisinage où elle est définie. Conclure seulement après cette transformation.

\end{enumerate}
\end{frame}
\begin{frame}{Continuité}
Pour \(r(x)=(x+1)/(x-2)\), le numérateur tend vers \(3>0\) en \(2\). À gauche, \(x-2<0\) et tend vers \(0\) : \(r(x)\to-\infty\). À droite, le dénominateur est positif : \(r(x)\to+\infty\). Écrire seulement « le dénominateur tend vers zéro » ne permettrait pas de choisir entre ces deux réponses.
\end{frame}
\begin{frame}{Continuité}
Pour \(x\ne0\),
\[
\dfrac{2x^3-x+1}{x^2+4}
=x\dfrac{2-1/x^2+1/x^3}{1+4/x^2}.
\]
Le quotient de droite tend vers \(2>0\) aux deux infinis; le facteur \(x\) impose \(+\infty\) à droite et \(-\infty\) à gauche. Diviser toujours par le degré du numérateur aurait aussi fonctionné, mais aurait laissé un dénominateur tendant vers zéro à analyser.
\end{frame}
\begin{frame}{Continuité}
Pour \(F(x)=x^3+x-1\), établir la continuité sur l’intervalle choisi et des signes opposés aux bornes. Pour l’unicité, si \(a<b\),
\[
F(b)-F(a)=(b-a)(a^2+ab+b^2+1)>0,
\]
car \(a^2+ab+b^2=(a+b/2)^2+3b^2/4\ge0\). Une fois l’unique zéro situé, tester le milieu d’un intervalle et conserver la moitié où les signes restent opposés. Cette dichotomie produit un encadrement, pas une valeur exacte par calculatrice.
\end{frame}
\begin{frame}{Continuité}
Une limite finie n’impose pas que la fonction soit définie au point. Si elle l’est, sa valeur peut différer de la limite. Pour rendre une fonction continue par prolongement, la valeur ajoutée doit être égale à cette limite; deux limites latérales différentes empêchent tout tel prolongement.
\end{frame}
\end{document}
